\section{Task 2 --- Evaluation of Effectiveness}
\label{sec:task2}

\subsection{Evaluation metrics (Task 2.1)}
\TODO{Definitions of the accuracy and beyond-accuracy metrics as implemented in \texttt{project/metrics}.}

\subsection{Comparison with baselines (Task 2.2)}
\TODO{Compare the best-performing models against the Random and Popularity baselines and discuss.}

% The table below is filled from report/tables/generated/model_results.tex once an experiment
% has produced it and it has been committed. Until then a red placeholder box is shown.
\begin{table}[htbp]
  \centering
  \caption{Accuracy and beyond-accuracy results per model.}
  \label{tab:model-results}
  \generatedtable{model_results}
\end{table}

\begin{figure}[htbp]
  \centering
  \generatedfigure[width=0.9\linewidth]{model_comparison}
  \caption{Model comparison.}
  \label{fig:model-comparison}
\end{figure}

\subsection{Hybrid coefficient analysis (Task 2.3)}
\TODO{How much does each individual recommender contribute to the hybrid's final prediction?}

\begin{figure}[htbp]
  \centering
  \generatedfigure[width=0.9\linewidth]{hybrid_coefficients}
  \caption{Learned hybrid coefficients.}
  \label{fig:hybrid-coefficients}
\end{figure}

\subsection{Further analyses of the hybrid and its components (Task 2.4)}
\TODO{Which additional analyses explain the behaviour of the hybrid and of its individual models?}

\subsection{Performance across user and item groups (Task 2.5)}
\TODO{For which user and item groups is each model more successful, in accuracy and beyond accuracy?}

\begin{figure}[htbp]
  \centering
  \generatedfigure[width=0.9\linewidth]{user_group_analysis}
  \caption{Performance per user group.}
  \label{fig:user-group-analysis}
\end{figure}

\subsection{Insights (Task 2.6)}
\TODO{What do these analyses suggest for novel or better hybrid models?}
